\begin{ex}
Trong phép nhân $2 \times 5 = 10$, số 10 được gọi là gì? Tính tích của 2 và 5.
\shortans{10}
\loigiai{
Tích của 2 và 5 là 10.
}
\end{ex}

\begin{ex}
Trong phép nhân $5 \times 4 = 20$, thừa số thứ nhất là 5, thừa số thứ hai bằng bao nhiêu?
\shortans{4}
\loigiai{
Thừa số thứ hai là 4.
}
\end{ex}

\begin{ex}
Biết hai thừa số là 2 và 7. Tính tích của hai số đó.
\shortans{14}
\loigiai{
Tích: $2 \times 7 = 14$.
}
\end{ex}

\begin{ex}
Biết hai thừa số là 5 và 6. Tính tích của hai số đó.
\shortans{30}
\loigiai{
Tích: $5 \times 6 = 30$.
}
\end{ex}

\begin{ex}
Biết hai thừa số là 2 và 9. Tính tích của hai số đó.
\shortans{18}
\loigiai{
Tích: $2 \times 9 = 18$.
}
\end{ex}

\begin{ex}
Biết hai thừa số là 5 và 8. Tính tích của hai số đó.
\shortans{40}
\loigiai{
Tích: $5 \times 8 = 40$.
}
\end{ex}

\begin{ex}
Trong phép nhân, nếu một thừa số là 0 thì tích bằng bao nhiêu?
\shortans{0}
\loigiai{
Bất kỳ số nào nhân với 0 cũng bằng 0.
}
\end{ex}

\begin{ex}
Trong phép nhân, nếu một thừa số là 1 thì tích bằng số nào? Cho phép nhân $7 \times 1$, kết quả bằng bao nhiêu?
\shortans{7}
\loigiai{
Ta có $7 \times 1 = 7$.
}
\end{ex}

\begin{ex}
Tính: $2 \times 3 + 4$.
\shortans{10}
\loigiai{
Ta có $2 \times 3 = 6$, rồi $6 + 4 = 10$.
}
\end{ex}

\begin{ex}
Tính: $5 \times 2 + 8$.
\shortans{18}
\loigiai{
Ta có $5 \times 2 = 10$, rồi $10 + 8 = 18$.
}
\end{ex}

\begin{ex}
Tính: $2 \times 8 - 6$.
\shortans{10}
\loigiai{
Ta có $2 \times 8 = 16$, rồi $16 - 6 = 10$.
}
\end{ex}

\begin{ex}
Tính: $5 \times 5 - 15$.
\shortans{10}
\loigiai{
Ta có $5 \times 5 = 25$, rồi $25 - 15 = 10$.
}
\end{ex}

\begin{ex}
Điền số thích hợp: $2 \times \dots = 12$.
\shortans{6}
\loigiai{
Số cần điền là: $12 : 2 = 6$.
}
\end{ex}

\begin{ex}
Điền số thích hợp: $5 \times \dots = 25$.
\shortans{5}
\loigiai{
Số cần điền là: $25 : 5 = 5$.
}
\end{ex}

\begin{ex}
Điền số thích hợp: $\dots \times 7 = 14$.
\shortans{2}
\loigiai{
Số cần điền là: $14 : 7 = 2$.
}
\end{ex}

\begin{ex}
Điền số thích hợp: $\dots \times 9 = 45$.
\shortans{5}
\loigiai{
Số cần điền là: $45 : 9 = 5$.
}
\end{ex}

\begin{ex}
Một nhóm có 2 bạn, 6 nhóm như thế có tất cả bao nhiêu bạn?
\shortans{12}
\loigiai{
Có: $2 \times 6 = 12$ bạn.
}
\end{ex}

\begin{ex}
Mỗi hộp có 5 chiếc bút bi, 7 hộp như thế có bao nhiêu chiếc bút bi?
\shortans{35}
\loigiai{
Có: $5 \times 7 = 35$ chiếc bút.
}
\end{ex}

\begin{ex}
Mỗi chuồng có 2 con thỏ, 8 chuồng như thế có tất cả bao nhiêu con thỏ?
\shortans{16}
\loigiai{
Có: $2 \times 8 = 16$ con thỏ.
}
\end{ex}

\begin{ex}
Bác An mua 5 túi gạo, mỗi túi nặng 5 kg. Bác An đã mua tất cả bao nhiêu ki-lô-gam gạo?
\shortans{25}
\loigiai{
Bác mua: $5 \times 5 = 25$ kg.
}
\end{ex}

\begin{ex}
Có 4 can nước, mỗi can chứa 5 lít. Hỏi cả 4 can chứa bao nhiêu lít nước?
\shortans{20}
\loigiai{
Chứa: $5 \times 4 = 20$ lít.
}
\end{ex}

\begin{ex}
Tìm tích của số bé nhất có hai chữ số và số 2.
\shortans{20}
\loigiai{
Số bé nhất có hai chữ số là 10. Tích: $10 \times 2 = 20$.
}
\end{ex}

\begin{ex}
Tìm tích của số liền sau số 1 và số liền trước số 6.
\shortans{10}
\loigiai{
Số liền sau 1 là 2, liền trước 6 là 5. Tích: $2 \times 5 = 10$.
}
\end{ex}

\begin{ex}
Tích của hai số là 18, biết một thừa số là 2. Thừa số kia là bao nhiêu?
\shortans{9}
\loigiai{
Thừa số kia là: $18 : 2 = 9$.
}
\end{ex}

\begin{ex}
Tích của hai số là 35, biết một thừa số là 5. Thừa số kia là bao nhiêu?
\shortans{7}
\loigiai{
Thừa số kia là: $35 : 5 = 7$.
}
\end{ex}
